\chapter{Architecture}

\section{High Level Structure}

The engine is split into a small set of modules under \filepath{src/demi}:

\begin{itemize}
  \item \filepath{assets}: asset manifest loading and reference extraction;
  \item \filepath{diagnostics}: structured diagnostics and severity reporting;
  \item \filepath{filesystem}: known source file discovery and file kind helpers;
  \item \filepath{schema}: validation of project, scene, HUD, save, and asset source data;
  \item \filepath{runtime}: scene loading, world data, rendering, physics, audio, Lua, and application loop;
  \item \filepath{cli}: command line entry point;
  \item \filepath{editor}: editor scaffold.
\end{itemize}

\section{Runtime Loop}

The runtime performs the following steps:

\begin{enumerate}
  \item load and validate the project file;
  \item load the main scene and optional HUD file;
  \item create the SDL window and renderer if SDL3 is available;
  \item load texture and audio assets from asset manifests;
  \item initialize Lua and load project, entity, and HUD button scripts;
  \item run fixed updates and physics at a fixed step;
  \item run Lua update callbacks;
  \item apply requested runtime changes such as window mode;
  \item render the world and HUD;
  \item shut down scripts, audio, renderer, and window.
\end{enumerate}

\section{World Model}

The runtime world is represented by \code{World} in \filepath{src/demi/runtime/SceneData.h}. It contains:

\begin{itemize}
  \item scene metadata;
  \item the HUD canvas size;
  \item entities and optional components;
  \item HUD rectangles, buttons, and text;
  \item debug lines.
\end{itemize}

Entities are simple value objects with optional components. The current component set includes transforms, cameras, sprites, Lua scripts, rigidbodies, box colliders, isometric grid data, hitbox controller data, and buildable metadata.

\section{Design Philosophy}

The current implementation prefers explicit, small APIs. New engine capabilities are expected to follow this pattern:

\begin{itemize}
  \item add source data fields first where possible;
  \item validate data through the CLI;
  \item expose only deliberate Lua APIs;
  \item keep generated outputs out of source directories;
  \item keep runtime/editor behavior backed by the same data and diagnostics paths.
\end{itemize}
